\section{Dataset Details}
\label{sec:app-dataset}

This section details the construction of FGMB. The benchmark is aggregated from a total of six data sources. These include four publicly available datasets-SAM, COYO, ImgDiff, and Vismin-and two proprietary datasets: XGoods and XNote, which we collected from an e-commerce platform to cover a broader range of fused-modal scenarios. Using the bounding box annotations from these datasets, we created regional query-candidate pairs. To guarantee that the benchmark is sufficiently challenging and discriminative, each query is paired with one to six hard negatives across all tasks.

\textbf{SAM.} This dataset is instrumental for building Tasks 1 and 2. We first converted the segmentation masks from SAM into bounding boxes by taking the maximum coordinates along each axis. Each region was then paired with a corresponding long caption from the DAM~\cite{lian2025dam} training set. To ensure an adequate supply of hard negatives, we required that each candidate pool include at least two distinct regions. For both Task 1 and Task 2, this process yielded 30k samples for training and 1.5k for evaluation.

\textbf{COYO.} From the COYO dataset, we utilized bounding boxes previously generated by FG-CLIP and subsequently re-captioned each region with DAM~\cite{lian2025dam}. Employing much longer captions allows for a richer alignment between the image region and its associated semantic concepts. These region-caption pairs naturally form the basis for text-to-image and image-to-text retrieval tasks, for which we prepared 1.5k evaluation samples each.

\textbf{ImgDiff.} As an image editing dataset, each sample in ImgDiff contains a source image, an editing instruction, an edited image, and a bounding box indicating the modified area. This structure lends itself naturally to an image-text-to-image retrieval task, where the query is a composite of the source image and the instruction, and the positive candidate contains the target region. When constructing the test split, we exclusively retained samples with two or more editing instructions to serve as hard negatives. This split comprises 21k training samples and 1.3k test samples.


\textbf{Vismin.} The structure of Vismin is analogous to that of ImgDiff; the primary difference is that Vismin contains real-world photographs, whereas ImgDiff's images are in the AI-generated style. We applied the same processing approach: the source image and edit instruction form the query, and the modified image serves as the positive candidate. From this dataset, we selected 46k samples for training and 1.3k queries for testing.

\textbf{XGoods.} Recognizing the limited availability of public datasets for region-level image-to-image retrieval, we collected XGoods from an e-commerce platform. Unlike datasets such as NIGHTS, which only contain the same object in similar backgrounds, XGoods offers more complex and realistic scenes. In this task, a commercial product image acts as the query, while a user-uploaded photo of the same item serves as the positive candidate. The bounding boxes of the query object were initially annotated automatically and then manually refined. We paired each query with six highly relevant products as hard negatives. For Task 3, we allocated 69k samples for training and 3.3k for evaluation. For Task 4, we also built a 3.4k data test split as a held-out dataset.


\textbf{XNote.} Similar to XGoods, XNote was also collected from social media content on the same e-commerce platform. The main difference is its focus on images from users' daily lives rather than commercial product displays. The data was processed using the same procedure as XGoods, from which we carefully selected 4.2k query-candidate pairs for the final test set.
% Each query region is paired with a candidate that serves as a detailed and unique description of the object inside the bounding box, or a visually similar object under different conditions (e.g., lighting or scene). 

% to create region pairs with consistent foregrounds but varied backgrounds, which are reserved as held-out data for evaluation.
% Queries are formed by combining the original image region with an associated instruction, and candidates include edited versions of the image. The model must reason over both the region and the background to succeed. 


\section{Training Details}
\label{sec:app-train}

\subsection{Training Parameters}
In the first stage, we use the batch size 64 and the learning rate of $1 \times 10^{-4}$. Only RAE's cross-attention modules and the projector are trainable.
In the second stage, the batch size is 1104 and the learning rate is $1.1 \times 10^{-4}$. We only update the language model backbone's parameters for one epoch using pure-text contrastive pairs. 
In the final stage, we fine-tune the language model backbone for two epochs on multimodal contrastive pairs, with a learning rate of $2 \times 10^{-4}$ and a batch size of 1104. 
To save the GPU memory for a larger batch size, which is crucial in contrastive learning, we use LoRA to fine-tune the LLM. The LoRA rank is set to be 64 and 128 for stage two and stage three, respectively. All three stages can be completed within 20 hours on 24 NVIDIA H20 GPUs. 

\subsection{Training Strategy Ablations}

During the training stage three, we develop the \textbf{mixed visual prompting strategy} to balance contextual and regional information. In particular, for tasks 1 to 3 in the FGMB training split, which rely more on local details, we provide only RAE tokens to the LLM. For task 4, which requires contextual understanding, we concatenate CE tokens with RAE tokens. This design allows the model to leverage CE tokens for context without injecting excessive background signals into RAE, thereby preserving fine-grained representation quality. 

We compare the mixed visual prompting strategy with some approaches in~\ref{tab:train-strategy}. This strategy proves to be able to provide consistent performance improvements across both meta tasks. As shown in Table~\ref{tab:train-strategy}, we compare several alternative training strategies. The ``crop'' setting uses only the RAE tokens for both meta tasks; the ``concat'' setting uses the concatenation of context encoder (CE) and RAE tokens as visual input; and the ``random'' strategy randomly selects either ``crop'' or ``concat'' with a 50\% probability for each meta task. We observe that our strategy achieves the best performance. This is mainly because of the trade-off between the two metatasks. Specifically, metatask 2 requires contextual information for comprehensive representation, while metatask 1 demands a highly localized, fine-grained understanding. Excessive context may thus act as a distraction in the latter. Therefore, we provide the CE tokens when training on metatask2, so that the model can access the global-level representations. Meanwhile, we use only RAE tokens on metatask 1 to prevent introducing background noise. %During evaluation, we use the same strategy.

\begin{table}[h]
\centering
\caption{Ablations on the training strategy in Stage-III using 3B models, evaluated on FGMB.
}
\label{tab:train-strategy}
\adjustbox{max width=\linewidth}{
\setlength{\tabcolsep}{3pt}
\begin{tabular}{lccccc}
\toprule
\textbf{Strategy} & \textbf{Task1} & \textbf{Task2} & \textbf{Task3} & \textbf{Task4} & \textbf{Avg.} \\
\midrule
crop  & 68.2 & 78.4 & 89.3 & 77.9 &  78.4 \\
concat  & 68.3 & 77.3 & 88.5 & 71.4 &  75.8  \\
random  & 67.8 & 78.9 & 89.2 & 73.0 &  76.7 \\
mixed visual prompts  & 71.0 & 77.4 & 90.0 & 80.8 &  79.9 \\
\bottomrule
\end{tabular}
}
\end{table}

